\documentclass[11pt, a4paper]{problems_template}

\begin{document}

\begin{titlepage}
  \begin{center}
	\Huge{Турнир городов}
  \end{center}
\end{titlepage}

\chapter{1980-1981. Номер 2}

\begin{exercise}
На бесконечной плоскости играют двое: один передвигает одну фишку-волка, другой - одну из $K$ фишек-овец. После хода волка ходит какая-нибудь из овец, затем, после следующего хода волка, опять какая-нибудь из овец и так далее. И волк, и овцы передвигаются за один ход в любую сторону не более чем на один метр. Верно ли, что для любого числа овец, участвующих в игре, существует такая первоначальная позиция, что волк не поймает ни одной овцы?
\end{exercise}


\chapter{1981-1982. Номер 3}

\begin{exercise}
Рассматривается последовательность
$$
1, \frac{1}{2}, \frac{1}{3}, \frac{1}{4}, \frac{1}{5}, \frac{1}{6}, \frac{1}{7}, \dots
$$
Существует ли арифметическая прогрессия
\begin{itemize}
\item длины 5
\item сколь угодно большой длины,составленная из членов этой последовательности?
\end{itemize}
\end{exercise}
 
\begin{Solution}
Существует для обоих пунктов. Возьмем дробь вида $\frac{1}{n!}$. Любая кратная ей дробь, имеющая вид
$\frac{k}{n!}$ будет членом последовательности если $k \leqslant n$. Таким образом получим арифметическую прогрессию длины $k$ с шагом $\frac{1}{n!}$.
\end{Solution}

\begin{exercise}
Квадрат разбит на $K^{2}$ равных квадратиков. Про некоторую ломаную известно, что она проходит через центры всех квадратиков (ломаная может пересекать сама себя). Каково минимальное число звеньев у этой ломаной?
\end{exercise}

\chapter{1982-1983. Номер 4}

\begin{exercise}
Числа от 1 до 1000 расставлены по окружности. Докажите, что их можно соединить 500 непересекающимися отрезками, разность чисел на концах которых (по модулю) не более 749.
\end{exercise}

\begin{Solution}
Распределим числа в два множества - от 251 до 750 в первом и все оставшиеся во втором. Пусть одно множество будет покрашено в зеленый цвет, а другое в синий. Заметим что мощности этих множеств равны и что разность любых двух чисел разного цвета удовлетворяет условию задачи. Очевидно, что найдутся два соседних числа разного цвета. Соединим их прямой. Выкинем эту пару из окружности. Действуя по аналогии соединим все числа непересекающимися отрезками. 
\end{Solution}

\chapter{1983-1984. Номер 5}

\begin{exercise}
Из листа клетчатой бумаги размером $29 \times 29$ клеточек вырезали 99 квадратиков $2 \times 2$ (режут по линиям). Докажите, что из оставшейся части листа можно вырезать еще хотя бы один такой же квадратик.
\end{exercise}

\chapter{1984-1985}

\begin{exercise}
Имеется 68 монет, причем известно, что любые две монеты различаются по весу. За 100 взвешиваний на \newline двухчашечных весах без гирь найдите самую тяжелую и самую легкую монеты.
\end{exercise}

\begin{Solution}
Взвесим 34 пары монет и получим две группы, которые назовем "Легкая" и "Тяжелая". Очевидно, что самая легкая монета находится в "Легкой" группе. Ее можно найти за 33 взвешивания. Аналогично за 33 взвешивания можно найти самую тяжелую монету в "Тяжелой" группе. Всего $34 + 33 + 33 = 100$ взвешиваний.
\end{Solution}

\chapter{1989-1990. Номер 11}

\begin{exercise}
Рассмотрим все возможные наборы чисел из множества $\{1, 2, 3, \dots, N\}$, не содержащие двух соседних чисел. Докажите, что сумма квадратов произведений чисел в этих наборах равна $(N+1)!-1$.
\end{exercise}

\begin{Solution}
Будем доказывать по индукции. База очевидна. Составим рекуррентное соотношение для искомой величины $S_{n}$. Она состоит из наборов которые включают $n$, самого $n$ и наборов, в которые $n$ не входит. Если $n$ входит в набор, то он учитывается в общем выражении как $n^{2}S_{n-2}$, а если нет, то используется $S_{n-1}$. В итоге
$$
S_{n} = n^{2}S_{n-2}+S_{n-1}+n^{2} = ((n-1)! - 1)n^{2} + (n! - 1) + n^{2} = n!(n + 1) - 1 = (n + 1)! - 1
$$  
\end{Solution}

\chapter{1998-1999. Номер 20}

\begin{exercise}
В круге провели $2n$ радиусов. Они разделили круг на $2n$ равных секторов: $n$ синих и $n$ красных, чередующихся в произвольном порядке. В синие сектора, начиная с некоторого, записывают против хода часовой стрелки числа от 1 до $n$. В красные сектора, начиная с некоторого, записывают те же числа, но по ходу часовой стрелки. Докажите, что найдется полукруг, в котором записаны все числа от 1 до $n$.
\end{exercise}

\begin{Solution}

\end{Solution}



\chapter{2000-2001. Номер 22}

\begin{exercise}
В весеннем туре турнира городов 2000 года старшеклассникам страны $N$ было предложено 6 задач. Каждую задачу решили ровно 1000 школьников, но никакие два школьника не решили вместе все 6 задач. Каково наименьшее возможное число старшеклассников страны $N$, принявших участие в весеннем туре?
\end{exercise}

\chapter{2020-2021}

\begin{exercise}
Существуют ли 100 таких натуральных чисел, среди которых нет одинаковых, что куб одного из них равен сумме кубов остальных?
\end{exercise}

\chapter{2024-2025}

\begin{exercise}
Учитель назвал две различные ненулевые цифры. Коля хочет составить делящееся на 7 семизначное число, в десятичной записи которого нет других цифр, кроме этих двух. Всегда ли Коля может это сделать, какие бы две цифры ни назвал учитель?
\end{exercise}



\end{document} 
 
 
